
We know that $ U$ is a random variable whose distribution is the uniform distribution on $ [0,1]$.
That is $ U$ is a mapping $ U: \Omega \to E$, with the following distribution 

\begin{align*}
    \P_{U} (A=(a,b)) = 
    \begin{cases}
	\int_{a}^{b} 1 dx & \mbox{if} \, a \geq 0 \, \mbox{and} \, b \leq 1\\
	0 & \mbox{else} 
    \end{cases}
\end{align*}
.
The quantile function or the inverse cumulative function is a mapping $ F^{-1}: (0,1) \to \R $,
meaning that we need to restrict the co-domain of $ U$ so that $ F^{-1}(U) $ is meaningful.
Thus, we have that now $ U: \Omega \to (0,1)$, and the distribution of $ \P_{U} $ is now never zero
since the we have restricted ourselves to the intervals $ (a,b)$ that are subsets of $ (0,1)$.

Like any other random variable $ F^{-1} (U)$ has a distribution that we can define. Let's find the
distribution of $ F^{-1} (U)$.

\begin{align*}
    \P_{F^{-1} (U)} (A) &= \P ( \{ \omega \in \Omega: F^{-1} (U(\omega) \in (a,b)) \} )\\
			&= \P( \{ \omega \in \Omega: U(\omega) \in (F(a),F(b)) \} 
\end{align*}
The last line follows from the fact that $ F^{-1} $ is a strictly increasing function so if
$ F^{-1} (U) > a$ then it implies that $ U > F(a)$.

We notice that $\P( U \in (F(a),F(b)) \}  $ is the same as the probability $ \P_{U} ((F(a), F(b))
$.

We get the following
\begin{align*}
    \P_{F^{-1}(U) } (A= (a,b)) &= \P_{U} ((F(a), F(b)))\\
			       &= \P_{U} ((e^{- e^{-a} }, e^{- e^{-b} }  ))\\
			       &= e^{-e^{-b} } - e^{-e^{-a} } \\
			       &= F(b) - F(a)
\end{align*}
.
The definition of a cumulative distribution function is implicitly assuming that there exists a
random number $ X$ such that $ X $ is a mapping $ X: \Omega \to E \subset \R$
\[
    F(x) = F_{X} (x) = \P_{X} ((-\infty, x])
.\] 
By Lemma 5.13, we have that

\[
    \P_{X} ((a,b)) = \P(X \in (a,b)) = F(b) -F(a) 
\] 
so we see that the probability distribution of $ F^{-1}(U)$ is the Gumbel distribution $ P_{X} $.
